
KV cache compression methods such as H2O eviction and quantization enable memory-efficient long-context inference in large language models. However, these methods apply uniform strategies across all tasks. This work investigates whether task structure should inform compression selection. Through gap statistic analysis of 21 LongBench tasks across 6 compression configurations, tasks cluster into three statistically distinct response groups (k* = 3, gap criterion satisfied: 0.957 >= 0.898). The clusters exhibit interpretable structure: multi-document QA and code tasks demonstrate high eviction sensitivity, while summarization tasks tolerate more aggressive compression. Additionally, first-100-token attention entropy discriminates task domains with large effect size (F = 38.05, p = 2.92 $\times 10^{-26}, \eta^2$ = 0.522), validating entropy as a potential routing signal. A third experiment testing whether high-entropy tasks tolerate eviction better yielded inconclusive results (p = 0.326) due to evaluation limitations. These findings establish an empirical foundation for task-conditioned KV cache compression: task-dependent structure exists and is detectable via attention features. The complete routing system remains future work.
